\documentclass{article}
\usepackage[utf8]{inputenc}
\usepackage{cite}
\usepackage{hyperref}

\title{Complex QA and Language Models Hybrid Architectures Survey}
\author{Your Name}
\date{\today}

\begin{document}

\maketitle

\begin{abstract}
This paper reviews the state-of-the-art of language models architectures and strategies for complex question-answering (CQA) with a focus on hybridization. Large Language Models (LLMs) like GPT-3 have demonstrated tremendous potential in general question-answering tasks but face significant challenges with complex, domain-specific questions. This survey explores required skills and evaluation techniques, and integrates findings from state-of-the-art projects such as BIG, BLOOM, and HELM. We discuss challenges like domain adaptation, efficient multi-step QA, long-form and non-factoid QA, safety and data protection, multimodal search, hallucinations, explainability, and truthfulness. Furthermore, we analyze current solutions and promising research trends including hybrid LLM architectural patterns, active human reinforcement learning, neuro-symbolic grounding, program synthesis, and iterated decomposition. 

\end{abstract}

\tableofcontents

\section{Introduction}
\label{sec:introduction}

Language models have revolutionized the field of natural language processing (NLP), with applications ranging from translation to text generation. However, when it comes to complex question-answering (CQA), these models face numerous challenges that necessitate more advanced architectures and strategies. Complex questions often require deep understanding, reasoning, and integration of specialized knowledge. Examples include questions on personal freedom across cultures or the optimal mix of power generation methods to mitigate climate change.

\subsection{State of Large Language Models (LLMs)}
The advent of large models like GPT-3 \cite{brown2020language} has shown the capacity of language models to handle a variety of question types. GPT-3, with its 175 billion parameters, can generate human-like text and answer questions based on the wide corpus it has been trained on. Nevertheless, its performance substantially declines when tasked with domain-specific, multi-step, or complex queries \cite{bubeck2023sparks}.

\subsection{Challenges in Complex Question-Answering (CQA)}
Complex QA involves many dimensions – it is not just about retrieving the correct fact, but also understanding nuanced questions that might span multiple domains and require sequential reasoning. These questions often need integration from different knowledge bases, understanding of temporal contexts, and the capability to handle ambiguity. Addressing these challenges requires innovative approaches beyond standard language models.

\section{Required Skills and Evaluation Techniques}
\label{sec:skills_evaluation}

\subsection{Required Skills for CQA}
For language models to be effective in CQA, they must possess several skills:
\begin{itemize}
    \item \textbf{Contextual Understanding}: The model must understand the context of the question to provide a relevant answer.
    \item \textbf{Reasoning and Inference}: The ability to perform multi-step reasoning and draw inferences is crucial.
    \item \textbf{Domain-Specific Knowledge}: Some questions require deep knowledge about specific fields.
    \item \textbf{Temporal Reasoning}: Understanding and reasoning with time-related information is essential.
    \item \textbf{Explainability and Transparency}: The model should be able to explain its reasoning to gain trust from users.
\end{itemize}

\subsection{Evaluation Techniques}
Evaluating complex QA models involves various metrics and benchmarks:
\begin{itemize}
    \item \textbf{Accuracy}: The correctness of the answers provided by the model.
    \item \textbf{Fairness}: Ensuring the model does not have biases in its responses \cite{bolukbasi2016man}.
    \item \textbf{Robustness}: The ability of the model to maintain performance under diverse and unseen scenarios \cite{hendrycks2020pretrained}.
    \item \textbf{Explainability}: The capability of the model to explain how an answer was derived.
    \item \textbf{Safety and Data Protection}: The model should handle sensitive data responsibly and ensure user privacy \cite{zhang2023model}.
\end{itemize}

\section{Benchmarking and Analysis of LLMs}
\label{sec:benchmarking}

\subsection{BIG, BLOOM, and HELM Projects}
Projects like BIG-Bench \cite{srivastava2022beyond}, BLOOM \cite{scao2022bloom}, and HELM \cite{bommasani2021opportunities} have created comprehensive benchmarks to analyze the performance of LLMs. These benchmarks aim to quantify the capabilities and limitations of language models in handling various complex tasks.

\subsection{Limits and Challenges in Tasks Complexity}
The analysis shows that while LLMs perform well on simple and standard tasks, their performance degrades on complex, multi-faceted tasks requiring deep reasoning and domain-specific knowledge \cite{bender2021dangers}. These limitations highlight the need for hybrid architectures that can leverage the strengths of different models and methodologies.

\section{Challenges in Complex QA}
\label{sec:challenges}

\subsection{Domain Adaptation}
One significant challenge in CQA is adapting to different domains. Standard models trained on broad data sets may lack the specificity required for niche fields. Techniques such as domain adaptation and transfer learning are being explored to address this \cite{gururangan2020don}.

\subsection{Decomposition and Efficient Multi-Step QA}
Complex questions often need to be broken down into simpler, more manageable parts. Decomposition strategies help in this regard but come with their own set of challenges in ensuring consistency and coherence among sub-answers \cite{min2021metaicl}.

\subsection{Long-Form and Non-Factoid QA}
Many complex queries require elaborate, detailed explanations rather than short factual answers. Handling long-form answers necessitates models that can generate coherent and contextually relevant paragraphs \cite{fan2019eli5}.

\subsection{Safety and Multi-Sensitivity Data Protection}
Given the sensitivity of some queries, it is crucial to ensure that models adhere to data protection standards and handle information securely \cite{sandven2021coping}.

\subsection{Multimodal Search}
Incorporating data from various modalities (text, image, audio, etc.) can provide more comprehensive answers. Multimodal models are being developed to integrate different types of data seamlessly \cite{alayrac2022flamingo}.

\subsection{Hallucinations, Explainability, and Truthfulness}
LLMs sometimes produce "hallucinations," or plausible-sounding but incorrect information. Ensuring that models are both accurate and capable of explaining their answers is a significant ongoing research area \cite{maynez2020faithfulness}.

\subsection{Temporal Reasoning}
Questions involving time-sensitive information pose additional layers of complexity. Models need to understand and reason about temporal contexts to provide accurate answers \cite{zhou2020liquid}.

\section{Current Solutions and Promising Research Trends}
\label{sec:solutions}

\subsection{Hybrid LLM Architectural Patterns}
Hybrid architectures that combine different types of models (e.g., rule-based systems, smaller specialized models) are gaining traction. These architectures can offer more robust and contextually accurate responses for complex queries \cite{biasse2022hybrid}.

\subsection{Training and Prompting Strategies}
Innovative training strategies, including few-shot, zero-shot learning, and advanced prompt engineering, can enhance the model's ability to handle complex questions effectively \cite{rae2021scaling}.

\subsection{Active Human Reinforcement Learning Supervised with AI}
Incorporating human feedback and reinforcement learning can help improve model accuracy and reliability. This hybrid approach leverages human expertise to guide AI learning processes \cite{christiano2017deep}.

\subsection{Neuro-Symbolic and Structured Knowledge Grounding}
Combining neural networks with symbolic reasoning systems can provide the best of both worlds: the flexibility of deep learning and the rigor of symbolic logic \cite{rosenblatt2021system}.

\subsection{Program Synthesis}
Automatically generating programs that can solve complex problems based on natural language instructions is another promising approach \cite{chen2021evaluating}.

\subsection{Iterated Decomposition}
Breaking down complex questions iteratively can improve the accuracy and coherence of answers. Techniques like MetaICL (Meta-learning for In-context Learning) have shown promise in this area \cite{yan2021metaicl}.

\section{Conclusion}
\label{sec:conclusion}

The landscape of complex question-answering (CQA) continues to evolve, with significant strides made in the development of large language models. However, addressing the intricate challenges of CQA requires more than extensions of existing technologies. Hybrid architectures, domain adaptation, advanced training, and prompt engineering, along with human supervision and incorporation of structured knowledge, represent promising avenues for future research. The state-of-the-art projects and methodologies discussed highlight the potential and the limitations present in current approaches, paving the way for more robust, accurate, and explainable systems in the near future.

\bibliographystyle{plain}
\bibliography{prompt1_try5_4o}

\end{document}